% !TeX root = ../main.tex

\setproblem{G}{\ProblemTitleG}
\section{\texorpdfstring{\pletter. \ptitle}{문제 \pletter}}

\begin{frame}[label=problem-G]
    \sectiontitle{\pletter}{\ptitle}
    \sectionmeta{
        출제자: \idtext{areian13}\\
        태그: \tagtext{\#math, \#geometry, \#rotating\_calipers, \#fft, \#polygon\_area}\\
        예상 난이도: \tiertext{Diamond4}
    }
    \begin{itemize}
        \item 제출 0번, 정답 0명 (정답률 0\%)
        \item 처음 푼 사람: -
    \end{itemize}
\end{frame}

\begin{frame}{\textbf{\pletter}. \ptitle}
    \textbf{정해 1 - Preliminaries: 순환 합성곱}
    \begin{itemize}
        \item 길이가 $N$인 두 배열 $a,b$의 순환 합성곱은 다음과 같습니다.
    
    \[
        (a*b)_i=\sum_{j=0}^{N-1}a_{(i-j)\bmod N}b_j
    \]
    
        \item 각 곱 $a_pb_q$는 $(p+q)\bmod N$번째 결과에 기여합니다.
        \item 모든 항의 쌍을 직접 계산하면 $O(N^2)$이 필요합니다.
    \end{itemize}
\end{frame}

\begin{frame}{\textbf{\pletter}. \ptitle}
    \textbf{정해 1 - 지지함수의 성질}
    \begin{itemize}
        \item 민코프스키 합의 지지값은 두 지지값의 합입니다.
    
    \[
        h_{A\oplus B}(u)
        =\max_{a\in A,\ b\in B}(a+b)\cdot u
        =h_A(u)+h_B(u)
    \]
    
        \item 음이 아닌 배율 $t$로 확대하면 지지값도 $t$배가 됩니다.
    
    \[
        h_{tA}(u)=t\,h_A(u)\qquad(t\ge0)
    \]
    
        \item 도형을 회전시키는 대신 방향을 반대로 회전시켜도 됩니다.
    \end{itemize}
    \[
        h_{R_j(A)}(u_i)=h_A(u_{(i-j)\bmod N})
    \]
\end{frame}

\begin{frame}{\textbf{\pletter}. \ptitle}
    \textbf{정해 1 - 만화경 곱에서 합성곱 찾기}
    \[
        \theta=\frac{2\pi}{N},\qquad
        a_i=h_P(u_i),\quad b_i=h_Q(u_i),\quad c_i=h_T(u_i)
    \]
    \begin{itemize}
        \item 원점이 $Q$의 내부에 있으므로 모든 $b_j$는 양수입니다.
    \end{itemize}
    \[
        c_i
        =\frac1N\sum_{j=0}^{N-1}b_jh_{R_j(P)}(u_i)
        =\frac1N\sum_{j=0}^{N-1}b_ja_{(i-j)\bmod N}
    \]
    \begin{itemize}
        \item 즉, 배열 $c$는 순환 합성곱과 완벽히 동일한 형태에, 마지막에 $N$으로 나눠 얻어집니다.
        \item 따라서 FFT를 사용하면 $O(N\log N)$에 계산할 수 있습니다.
    \end{itemize}
\end{frame}

\begin{frame}{\textbf{\pletter}. \ptitle}
    \textbf{정해 1 - 지지값 배열 구하기}
    \begin{itemize}
        \item 방향이 반시계로 회전하면 내적이 최대인 꼭짓점도 같은 방향으로 이동합니다.
        \item 회전하는 캘리퍼스로 각 다각형을 한 바퀴 순회하며 모든 지지값을 구합니다.
    
    \[
        \boxed{O(N+V+W)}
    \]
    \end{itemize}
\end{frame}

\begin{frame}{\textbf{\pletter}. \ptitle}
    \textbf{정해 1 - 지지직선과 변의 길이}
    \begin{itemize}
        \setlength{\abovedisplayskip}{5pt}
        \setlength{\belowdisplayskip}{5pt}
        \item $C=P\circledast_N Q$의 지지직선을 $L_i:x\cdot u_i=c_i$라 합시다.
        \item $L_i$ 위에서 원점의 수선의 발을 기준으로 좌표를 잡습니다.
        \item $L_{i-1},L_{i+1}$과의 교점 좌표는 각각 다음과 같습니다.
    
    \[
        \frac{c_i\cos\theta-c_{i-1}}{\sin\theta},
        \qquad
        \frac{c_{i+1}-c_i\cos\theta}{\sin\theta}
    \]
    
        \item 두 교점 사이의 길이를 $\ell_i$라고 하면
    
    \[
        \ell_i
        =\frac{c_{i-1}+c_{i+1}-2\cos\theta\,c_i}{\sin\theta}
    \]
    
        \item 모든 $\ell_i>0$이면 이 교점들은 볼록다각형의 경계를 이룹니다.
    \end{itemize}
\end{frame}

\begin{frame}{\textbf{\pletter}. \ptitle}
    \textbf{정해 1 - 변의 길이가 양수인 이유 (1)}
    \begin{itemize}
        \setlength{\abovedisplayskip}{5pt}
        \setlength{\belowdisplayskip}{5pt}
        \item $u_i$ 방향에서 $P$의 지지값을 결정하는 꼭짓점을 $p$라 합시다.
    
    \[
        a_{i-1}\ge p\cdot u_{i-1},\qquad
        a_{i+1}\ge p\cdot u_{i+1}
    \]
    
        \item $u_{i-1}+u_{i+1}=2\cos\theta\,u_i$이므로
    
    \[
        e_i:=a_{i-1}+a_{i+1}-2\cos\theta\,a_i\ge0.
    \]
    
        \item 원점이 $P$의 내부에 있으므로 모든 $a_i$는 양수입니다.
    
    \[
        \sum_i e_i=2(1-\cos\theta)\sum_i a_i>0
    \]
    
    \end{itemize}
\end{frame}

\begin{frame}{\textbf{\pletter}. \ptitle}
    \textbf{정해 1 - 변의 길이가 양수인 이유 (2)}
    \begin{itemize}
        \item $c=\frac1N(a*b)$에 앞의 관계를 적용합니다.
    
    \[
        c_{i-1}+c_{i+1}-2\cos\theta\,c_i
        =\frac1N\sum_{j=0}^{N-1}e_{i-j}b_j>0
    \]
    
        \item 적어도 하나의 $e_i>0$이고 모든 $b_j>0$이므로 위 합은 양수입니다.
        \item 따라서 모든 $\ell_i>0$이며, 인접한 교점들을 이으면 됩니다.
        \item 변 방향이 차례로 회전하므로 이 다각형은 볼록입니다.
        \item 이 다각형이 주어진 반평면들의 교집합 $C$입니다.
    \end{itemize}
\end{frame}

\begin{frame}{\textbf{\pletter}. \ptitle}
    \textbf{정해 1 - 넓이 계산}
    \begin{itemize}
        \setlength{\abovedisplayskip}{5pt}
        \setlength{\belowdisplayskip}{5pt}
        \item 원점은 $C$의 내부에 있고, 각 변까지의 거리는 $c_i$입니다.
        \item 원점과 각 변이 만드는 삼각형의 넓이를 더합니다.
    
    \[
        S=\frac12\sum_i c_i\ell_i
    \]
    
        \item 변의 길이를 대입하고 순환 인덱스를 정리하면
    
    \[
        \boxed{
            S=\frac{
                \displaystyle\sum_i c_ic_{i+1}
                -\cos\theta\displaystyle\sum_i c_i^2
            }{\sin\theta}
        }
    \]
    
        \item 지지값 배열 $c$만 있으면 $O(N)$에 계산할 수 있습니다.
    \end{itemize}
\end{frame}

\begin{frame}{\textbf{\pletter}. \ptitle}
    \textbf{정해 1 - 수치적으로 안정적인 넓이 계산}
    \begin{itemize}
        \setlength{\abovedisplayskip}{5pt}
        \setlength{\belowdisplayskip}{5pt}
        \item $N$이 크면 $\cos\theta$가 $1$에 가까워집니다.
        \item 원래 식은 큰 두 합을 빼면서 소거 오차가 커질 수 있습니다.
        \item 다음 항등식과 반각 공식을 사용합니다.
    \end{itemize}
    \[
        \sum_i(c_i-c_{i+1})^2
        =2\sum_i c_i^2-2\sum_i c_ic_{i+1}
    \]
    \[
        \frac{1-\cos\theta}{\sin\theta}=\tan\frac\theta2
    \]
    \[
        \boxed{
            S=\tan\frac\theta2\sum_i c_i^2
            -\frac1{2\sin\theta}\sum_i(c_i-c_{i+1})^2
        }
    \]
\end{frame}

\begin{frame}{\textbf{\pletter}. \ptitle}
    \textbf{정해 1 - 전체 알고리즘}
    \begin{enumerate}
        \item 회전하는 캘리퍼스로 두 지지값 배열을 구합니다.
        \item FFT로 순환 합성곱을 계산하고 $N$으로 나눕니다.
        \item 구한 지지값으로 반평면 교집합의 넓이를 계산합니다.
    \end{enumerate}
    \begin{itemize}
        \item 시간 복잡도와 공간 복잡도는 각각 다음과 같습니다.
    \end{itemize}
    \[
        \boxed{O(V+W+N\log N)},\qquad
        \boxed{O(V+W+N)}
    \]
\end{frame}

\begin{frame}{\textbf{\pletter}. \ptitle}
    \textbf{정해 2 - 지지값 배열의 기하적 성질}
    \begin{itemize}
        \item 지지값 배열 $a,b,c$와 각도 $\theta$는 정해 1의 정의와 동일합니다.
        \item 지지점이 $p=(x,y)$일 때, 지지값은 다음과 같습니다.
        \[
            a_i=x\cos(i\theta)+y\sin(i\theta).
        \]
        \item 세 방향 $i-1,i,i+1$의 지지점이 같으면, 덧셈 정리로 다음 등식을 얻습니다.
        \[
            a_{i-1}+a_{i+1}=2\cos\theta\,a_i.
        \]
        \item 이 관계를 이용해 지지점이 바뀌는 위치만 남길 수 있습니다.
    \end{itemize}
\end{frame}

\begin{frame}{\textbf{\pletter}. \ptitle}
    \textbf{정해 2 - 지지값 배열의 희소화}
    
    \begin{itemize}
        \item 배열 $x$에 대해 아래 배열을 정의해봅시다. 
        \[
        (Dx)_i:=x_{i-1}-2\cos\theta\,x_i+x_{i+1}
    \]
        \item 지지값 배열 $a$에 대해 세 방향의 지지점이 같으면 $(Da)_i=0$입니다.
        \item 지지점의 변경은 최대 $V$번이고, 각 경계는 두 항에만 영향을 줍니다.
        \item $Da,Db$는 길이가 $N$이지만, 실제로는 $0$이 아닌 항의 인덱스와 값만 남겨 압축합니다.
        \item 이때 저장하는 항의 수는 각각 $2V$개, $2W$개 이하가 됩니다. 
    \end{itemize}
\end{frame}

\begin{frame}{\textbf{\pletter}. \ptitle}
    \textbf{정해 2 - 희소 배열의 합성곱}
    \begin{itemize}
        \item $D$도 고정된 배열과의 순환 합성곱이므로 다음 관계가 성립합니다.
        \[
            \boxed{(Da)*(Db)=D^2(a*b)=N\,D^2c}
        \]
        \item 이 합성곱을 $N$으로 나눈 배열을 $f:=D^2c$로 정의합니다.
        \item $0$이 아닌 항의 쌍만 곱하므로 $f$를 $O(VW)$에 구할 수 있습니다.
        \item 이제 $f$로부터 원래의 지지값 배열 $c$를 복원하면 됩니다.
    \end{itemize}
\end{frame}

\begin{frame}{\textbf{\pletter}. \ptitle}
    \textbf{정해 2 - 점화식으로 지지값 복원하기}
    \begin{itemize}
        \item $y=Dc$라 놓으면 $Dy=f$이며, 각 식을 다음 항에 대해 정리하면
        \[
            \begin{aligned}
                y_{i+1}&=2\cos\theta\,y_i-y_{i-1}+f_i,\\
                c_{i+1}&=2\cos\theta\,c_i-c_{i-1}+y_i.
            \end{aligned}
        \]
        \item 초깃값 $y_0,y_1,c_0,c_1$을 알면 나머지 항을 순서대로 계산합니다.
        \item 먼저 배열 $y$를 복원한 뒤 $c$를 복원하며, 총 $O(N)$이 걸립니다.
        \item 복원한 $c$로 넓이를 구하는 과정은 정해 1과 같습니다.
    \end{itemize}
\end{frame}

\begin{frame}{\textbf{\pletter}. \ptitle}
    \textbf{정해 2 - 초깃값 직접 계산하기}
    \begin{itemize}
        \item 복원에 필요한 네 초깃값은 합성곱의 정의로 직접 계산합니다.
        \[
            \begin{aligned}
                c_t&=\frac1N\sum_{j=0}^{N-1}a_jb_{t-j},\\
                y_t&=\frac1N\sum_{j=0}^{N-1}(Da)_jb_{t-j}
                \qquad(t=0,1).
            \end{aligned}
        \]
        \item $c_0,c_1$은 $O(N)$, $y_0,y_1$은 $Da$의 희소성을 이용해 $O(V)$입니다.
        \item 필요한 항만 직접 계산하므로 초깃값을 구하는 데는 총 $O(N+V)$입니다.
    \end{itemize}
\end{frame}

\begin{frame}{\textbf{\pletter}. \ptitle}
    \textbf{정해 2 - 수치 오차를 줄이는 희소화}
    \begin{itemize}
        \item $u_i$ 방향의 지지점을 $p_i$라 하면, $Da$를 다음과 같이 계산합니다.
        \[
            (Da)_i
            =(p_{i-1}-p_i)\cdot u_{i-1}
            +(p_{i+1}-p_i)\cdot u_{i+1}
        \]
        \item $u_{i-1}+u_{i+1}=2\cos\theta\,u_i$를 이용해 변형한 동일한 식입니다.
        \item 큰 지지값 세 개를 직접 빼는 계산을 줄일 수 있습니다.
        \item 세 지지점이 모두 같으면 계산 없이 $0$으로 처리하며, $Db$도 같습니다.
    \end{itemize}
\end{frame}

\begin{frame}{\textbf{\pletter}. \ptitle}
    \textbf{정해 2 - 수치 오차를 줄이는 복원}
    \begin{itemize}
        \item $Dx=g$의 복원에서는 차이 $\delta_i=x_i-x_{i-1}$를 유지합니다.
        \[
            \alpha=4\sin^2\frac\theta2=2-2\cos\theta
        \]
        \[
            \begin{aligned}
                \delta_{i+1}&=\delta_i+g_i-\alpha x_i,\\
                x_{i+1}&=x_i+\delta_{i+1}.
            \end{aligned}
        \]
        \item $\alpha$는 반각 공식으로 구하고, 복원은 차이를 갱신하며 계산합니다.
        \item 이 점화식을 $y$와 $c$의 복원에 각각 사용합니다.
    \end{itemize}
\end{frame}

\begin{frame}{\textbf{\pletter}. \ptitle}
    \textbf{정해 2 - 전체 알고리즘}
    \begin{enumerate}
        \item 정해 1과 같이 지지값 배열과 각 방향의 지지점을 구합니다.
        \item 희소 배열 $Da,Db$를 만들고, 항의 쌍을 직접 곱해 $f=D^2c$를 구합니다.
        \item 초깃값 네 개를 구하고, 점화식으로 배열 $y,c$를 차례로 복원합니다.
        \item 정해 1의 넓이 공식을 적용합니다.
    \end{enumerate}
    \begin{itemize}
        \item 시간 복잡도와 공간 복잡도는 각각 다음과 같습니다.
    \end{itemize}
    \[
        \boxed{O(N+VW)},\qquad
        \boxed{O(N+V+W)}
    \]
\end{frame}

\begin{frame}{\textbf{\pletter}. \ptitle}
    \textbf{출제자의 코멘트}
    \begin{itemize}
        \item $N$이 굳이 2의 거듭제곱 꼴이 아니어도 수학적으로 성립하는 문제입니다만,\\
        부차적인 순환 쪽 구현이 크게 의미 있진 않다 생각해 제한하였습니다.
        \item 또한 지지값이 모두 0 이상인 입력에 대해 성립하고, \\
        `어떤 점 집합의 컨벡스 헐이 원점을 포함한다면 항상 모든 지지값은 0 이상'이므로\\
        이 문제를 점 집합에 대해 일반화할 수 있습니다만,\\
        그냥 `볼록 다각형 $\times$ 볼록 다각형 = 볼록 다각형'인게 예쁘기도 하고, \\
        위 조건을 입력에 주저리 주저리 쓰고 컨벡스헐 돌리게 만드는게 별로여서\\
        원점을 포함한 컨벡스 폴리곤으로 제한했습니다.
        \item 거듭 만화경 $n$제곱 $P_{\circledast_N}^n$을 정의할 수 있겠죠? \\
        이때 $\lim_{n\to\infty}P_{\circledast_N}^n$의 모양은 정$N$각형으로 수렴한다는 특징이 있습니다. 
    \end{itemize}
\end{frame}